\chapter{Репродуцибилност, придонес и заклучок}
\label{chap:conclusion}

\section{Клонирање и pinned верзии}

Референтниот snapshot се добива со:

\begin{lstlisting}[language=bash,caption={Клонирање на јавната feature гранка со submodules},label={lst:clone}]
git clone --recurse-submodules \
  --branch feature/emotion-robot-integration \
  https://github.com/dimitarbez/lite3-ros-sim.git ROS
cd ROS
git rev-parse HEAD
git submodule status
\end{lstlisting}

Очекуваните SHA-вредности се wrapper
\code{c41f26ef4414e5c9a474}\allowbreak\code{8f3c88aee33820526c01}, EmotionBot
\code{f267d28552bbf46f30a2}\allowbreak\code{63bb291e71ab91a2b91b} и Lite3\_VMC
\code{785a3846ec2c44a03806}\allowbreak\code{18490f837786b3250f76}. Тие се проверуваат пред
build, бидејќи branch head може подоцна да се помести. \code{git status} се
проверува посебно во wrapper-от и во секој nested repository. Не се извршува
\code{bootstrap} врз неразгледана dirty/detached работа.

\section{Build и симулација}

Од workspace root:

\begin{lstlisting}[language=bash,caption={Репродуцибилен setup и offline simulation}]
make -C lite3-noetic build
make -C lite3-noetic start
make -C lite3-noetic bootstrap

# Терминал 1: цел интегриран offline graph
make -C lite3-noetic run-emotion-sim

# Терминал 2: интерактивен разговор
make -C lite3-noetic run-emotion-chat
\end{lstlisting}

За live OpenAI се користи \code{run-emotion-sim-openai}; key се става во
ignored, mode-600 \code{lite3-noetic/ws/Lite3\_VMC/.env}, никогаш во трудот
или Git. Пред chat handoff се проверуваат четирите status flags. Ако watchdog
ја одземе motion permission, се чека recovery и се користи service-от
\code{/emotion\_bot/set\_motion\_enabled}; readiness не се bypass-ира.

Безбедно завршување значи disable, Ctrl-C на chat/simulation и финална нулта
команда. Ако ROS graph е stale или има XML-RPC refusal, се користи еден чист
\code{make -C lite3-noetic restart}, а не паралелни graphs.

\section{Verification команди}

Главните offline gates се:

\begin{lstlisting}[language=bash,caption={Repeatable verification targets}]
make -C lite3-noetic verify-emotion
make -C lite3-noetic verify-hardware-offline
\end{lstlisting}

Првиот target го проверува build-от, static/unit/process-boundary/ROS tests и
двата Gazebo sweeps. Вториот го гради одделниот hardware workspace и ги
проверува packet, mapping, contact, profile, ownership, release и watchdog
механизмите без robot interfaces. \textbf{Ниту еден target не е доказ за
physical commissioning.}

Овој труд намерно не дава copy-paste процедура за физичко активирање. Таква
операција бара одделна тековна авторизација, свеж преглед на hardware runbook и
manuals, exact robot state, charged battery, clear area, physical restraint,
STOP/operator readiness и потврда дека нема друг sender. Датираните команди во
project docs се evidence records, а не standing permission.

\section{Оригинален придонес и атрибуција}

Разграничувањето во \cref{tab:contributions} е важно и авторски и лиценцно.
Не е коректно upstream Lite3 controller, Gazebo model или EmotionBot logic да
се припишат на авторот на интеграцијата. Оригиналниот придонес е во
архитектурното поврзување, contracts, safety layers, test infrastructure,
hardware runner, профилите и evidence discipline.

\begin{table}[htbp]
\centering
\caption{Разграничување на upstream основата и проектниот придонес}
\label{tab:contributions}
\small
\begin{tabularx}{\textwidth}{L{0.28\textwidth}YY}
\toprule
\textbf{Компонента} & \textbf{Потекло} & \textbf{Статус во проектот} \\
\midrule
Lite3 Gazebo модел, основни controllers и keyboard & DEEP Robotics Lite3\_VMC & Преземена upstream основа; зачувана baseline споредба. \\
Emotion appraisal, personality и memory & EmotionBot, GPL-3.0 & Преземена domain основа; додаден reusable headless EmotionEngine. \\
ROS Noetic Docker/WSL workflow & проектен придонес & Wrapper images, lifecycle, Make targets и docs. \\
\code{emotion\_bot\_ros} & проектен придонес & Contracts, adapters, chat coordinator, mapper, safety bridge и tests. \\
OpenAI sidecar & проектен придонес & Python 3.12 process boundary, stream/fallback/cancellation и secret isolation. \\
Controller expression actions & проектна измена во fork & Bounded hop/stomp, cancellation/recovery и C++ tests. \\
Split-host hardware workspace & проектен придонес & Passive telemetry, STOP relay, shared state, services и fail-closed diagnostic graph. \\
MotionSDK expression runner & проектен придонес врз official SDK & Exclusive owner, stand sequence, profiles, contact/feedback gates и status. \\
Physical profiles/tickets & проектен придонес & Neutral/Joy/Sadness/Fear/Anger acceptance и четири planned profiles. \\
Verification records & проектен придонес & Dated evidence, failure history и reproducibility boundary. \\
\bottomrule
\end{tabularx}
\end{table}


EmotionBot GPL-3.0 notices мора да се зачуваат; пред копирање или vendoring се
проверуваат redistribution implications. DEEP Robotics MotionSDK и Lite3\_VMC
се користат како официјална upstream основа. API и simulator documentation се
проверуваат во официјалните извори, а не во непроверени tutorials.

\section{Идна работа}

Најнепосредни задачи се:

\begin{enumerate}
  \item bounded physical revalidation на Joy и Anger persistent contact-miss
  recovery, без промена на thresholds;
  \item consolidated operator verdict за normal chat sweep и дури потоа
  одлука дали checked-in allowlist да се прошири;
  \item имплементација, native/aarch64 тестови, bounded physical telemetry и
  visual acceptance за Affection, Curiosity, Disgust и Surprise;
  \item формализирана multi-operator visual study со однапред дефинирани
  критериуми, наместо само commissioning operator verdict;
  \item независна contact sensing валидација на torque-derived estimator;
  \item миграција од Gazebo Classic кон поддржан simulator без губење на
  controller/telemetry acceptance tests; и
  \item понатамошно докажување на STOP latency, battery policy и release
  postures во точно дефинирани, restrained сесии.
\end{enumerate}

Airborne hop, locomotion, yaw и dynamic physical actions не се логичен „следен
мал чекор“; тие се одделни commissioning проекти со нови flight/landing,
balance и operator-risk барања.

\section{Заклучок}

Во трудот е прикажан целосен пат од разговорен input до симулиран и физички
емоционален израз, со јасна граница меѓу reasoning и actuation. EmotionEngine
обезбедува валентност, активација и девет категории. Versioned ROS contracts,
turn ordering и OpenAI sidecar овозможуваат inspectable разговор без
dependency/secret компромис во Noetic. Mapper-от и safety bridge-от ги
претвораат категориите во девет центрирани Gazebo профили, а автоматизираните
tests ги покриваат movement, transition, watchdog, zeroing и no-hardware
boundary.

Физичкиот систем користи поинаков, split-host пат со единствен MotionSDK owner,
state/feedback/STOP/battery/attitude/tracking gates и torque-derived contact
estimate. Neutral, Joy, Sadness, Fear и Anger имаат физички извршени и
операторски прифатени реакции. Четирите преостанати категории искрено се
прикажани како Neutral fallback. Последните Joy/Anger recovery корекции не се
прогласени за физички потврдени само затоа што се deployed и offline-tested.

Најважниот инженерски резултат не е само дека роботот се движи. Резултат е
доказната дисциплина со која се разликува command од motion, motion од
recognizable emotion и implementation од commissioning. Неуспешните height,
takeover, contact, attitude, battery, DNS и clock експерименти директно го
обликувале финалниот дизајн. Поради тоа проектот претставува сериозен,
репродуцибилен и безбедносно ограничен проектен прототип, со јасно
дефинирани граници за идно проширување.
